\begin{lstlisting}[language=python]

import hashlib
import unicodedata

from cryptography.hazmat.primitives.ciphers.aead import AESGCM
from sympy import primerange


def ejercicio_cuatro() -> None:

    try:
        with open("file.txt", "r", encoding="utf-8") as f:
            texto = f.read()
        with open("hashes.txt", "r", encoding="utf-8") as f:
            # Cargamos los hashes validos en un set para busqueda eficiente.
            # Se asume que cada linea del archivo contiene un hash SHA-256 valido.
            hashes_validos = set(line.strip() for line in f if line.strip())
        # Cargamos los datos cifrados en formato hexadecimal, eliminando cualquier espacio o salto de linea.
        with open("cipher.txt", "r", encoding="utf-8") as f:
            cipher_hex = f.read().replace("\n", "").strip()
    except FileNotFoundError as e:
        print(f"[-] Error: Archivo no encontrado - {e}")
        return

    # Se eliminan acentos, simbolos y saltos de linea.
    # Conservamos unicamente letras minusculas (a-z) para asegurar un bloque continuo.
    texto_norm = (unicodedata.normalize("NFKD", texto).encode("ascii", "ignore").decode("utf-8"))
    texto_filtrado = "".join(c for c in texto_norm if c.isalpha()).lower()
    
    # Generamos los primeros 50 numeros primos utilizando sympy.primerange
    primos = [p for _, p in zip(range(50), primerange(2, 1000))]
    palabra_base = "kevinmitnick"

    # El uso de `% len(texto_filtrado)` garantiza que el texto funcione como un "loop infinito".
    # por que si el indice calculado excede la longitud del texto, simplemente vuelve al inicio.
    opciones_extraccion = [
        # A: El indice extraido es el numero primo restandole 2.
        "".join(texto_filtrado[(p - 2) % len(texto_filtrado)] for p in primos),
        # B: Se usa el numero primo como indice, y se aplica un desplazamiento Cesar de (p-2).
        "".join(
            chr((ord(texto_filtrado[p % len(texto_filtrado)]) - 97 + (p - 2)) % 26 + 97)
            for p in primos
        ),
        # C: Indice en (p-2) y desplazamiento Cesar en (p-2).
        "".join(chr((ord(texto_filtrado[(p - 2) % len(texto_filtrado)]) - 97 + (p - 2)) % 26+ 97)
            for p in primos
        ),
    ]

    clave_correcta = None
    hash_correcto = None

    for extraido in opciones_extraccion:
        # Exploramos las diferentes formas de "mezclar" el texto con la palabra base "kevinmitnick"
        mezclas = [
            extraido,
            # Mezcla por suma (Estilo Cifrado Vigenere). Por ejemplo, si extraido = "abcde..." y
            # palabra_base = "kevinmitnick", la mezcla seria: "kkbcvindemnitck..." 
            # (alternando caracteres de ambos strings pero sumando sus valores).
            "".join(chr((ord(extraido[i])- 97+ ord(palabra_base[i % len(palabra_base)])- 97)% 26+ 97)
                for i in range(50)),
            # Mezcla por resta. Por ejemplo, si extraido = "abcde..." y palabra_base = "kevinmitnick", la mezcla seria:
            # "akbcvindemnitck..." (alternando caracteres de ambos strings pero restando sus valores).
            "".join(chr((ord(extraido[i]) - ord(palabra_base[i % len(palabra_base)])) % 26+ 97)
                for i in range(50)),
            # Se toma el caracter i-esimo de 'extraido' y se concatena con el caracter i-esimo de 'palabra_base' (ciclico).
            # Por ejemplo, si extraido = "abcde..." y palabra_base = "kevinmitnick", la mezcla seria:
            # "akebcvindemnitck..." (alternando caracteres de ambos strings)
            "".join(extraido[i] + palabra_base[i % len(palabra_base)] for i in range(25)),
        ]

        # Comparamos el SHA-256 de cada combinacion generada contra la lista de control
        for candidata in mezclas:
            h = hashlib.sha256(candidata.encode()).hexdigest()
            if h in hashes_validos:
                clave_correcta = candidata
                hash_correcto = h
                break
        if clave_correcta:
            break

    if not clave_correcta:
        print("[-] No se encontro ninguna clave que coincida con los hashes.")
        return

    print(f"\n[+] Clave confirmada: {clave_correcta}")
    print(f"[+] Hash coincidente : {hash_correcto}")

    try:
        data = bytes.fromhex(cipher_hex)
    except ValueError:
        print("[-] Error: cipher.txt no contiene un formato hexadecimal valido.")
        return

    nonce = data[:12]  # AES-GCM estandar usa un Nonce (IV) de 12 bytes
    ciphertext = data[12:]  # El resto del bloque es el texto cifrado + Tag de autenticacion

    llaves_a_probar = {
        "Hash Directo a Bytes (Lectura Literal)": bytes.fromhex(hash_correcto),
        "Doble Hash (Lectura sugerida por Hint)": hashlib.sha256(hash_correcto.encode()).digest(),
        "Derivacion Original (Hash de Clave Extr)": hashlib.sha256(clave_correcta.encode()).digest(),
    }

    exito = False
    for _, key_bytes in llaves_a_probar.items():
        aesgcm = AESGCM(key_bytes)
        try:
            # Si decrypt() no lanza la excepcion 'InvalidTag', significa que la llave derivo exitosamente.
            decrypted = aesgcm.decrypt(nonce, ciphertext, None)
            print("[+] Mensaje Descifrado:\n")
            print(f"{decrypted.decode('utf-8')}")
            exito = True
            break
            # La llave es incorrecta para este ciphertext. Ignoramos y probamos la siguiente.
        except Exception:
            continue

    if not exito:
        print("[-] Error: Ninguna de las llaves derivadas logro descifrar el mensaje.")

if __name__ == "__main__":
    ejercicio_cuatro()

\end{lstlisting}